% Volume III: Attention and Publicity
% Part V. Publicity Without Comprehension
% Chapter 25: Attention as a Scarce Public Resource
% Spine: proof-spines/ch025-spine.md

\chapter{Attention as a Scarce Public Resource}
\label{ch:ch025}

Public attention is modeled as a finite \textbf{attention allocation vector}
\[
\mathbf a=(a_1,\ldots,a_n),
\qquad
\sum_{i=1}^{n} a_i \leq A.
\]
\Definition\ Publicity does not abolish scarcity; it organizes competition for \(A\). Because every increase in salience somewhere spends capacity unavailable elsewhere, agenda-setting, repetition, and timing become epistemic mechanisms rather than merely rhetorical ones.

Powerful actors may therefore dominate not only by suppressing information, but by capturing or diluting attention shares within the common budget \(A\). Equal formal access to speak does not equal equal access to uptake. The chapter reframes attention as public infrastructure and motivates triage, watchdog specialization, and anti-monopolization norms for the collective sensorium.
